\subsection{Quaternions e justificativa de sua utilização}
\label{sec:representacao_quaternion}
 
Os quaternions constituem uma alternativa para representar a orientação, pois evitam as descontinuidades associadas às singularidades dos ângulos de Euler e permitem descrever a atitude do perseguidor.

Um quaternion unitário é definido como um elemento do espaço quaternional de quatro dimensões, pertencente ao conjunto
dos quaternions unitários $q = [w,\vec{v}\,]$, com $\|q\| = 1$, em que $w$ é a parte escalar e $\vec{v}$ é a parte vetorial. Cada quaternion unitário $q$ corresponde a uma rotação de ângulo $\theta$ em torno de um eixo unitário $\vec{u}$, dada por:

\begin{equation}
q(\theta,\vec{u}) = \cos\!\left(\frac{\theta}{2}\right) + \sin\!\left(\frac{\theta}{2}\right)\,\vec{u}.
\label{eq_quat_axis_angle}
\end{equation}

Essa parametrização fornece uma representação livre de singularidades geométricas para a integração contínua das equações diferenciais de atitude.
Em aplicações numéricas, impõe-se a condição de unidade $\|q\|=1$ e utiliza-se o desvio de norma como critério de monitoramento de erros de integração ao longo da simulação.

Do ponto de vista físico, os quaternions preservam a norma e realizam rotações por multiplicação, assegurando uma transição suave entre orientações sucessivas.
Embora exijam uma variável adicional em comparação aos ângulos de Euler, sua principal vantagem reside na ausência de singularidades e no fato de que as equações cinemáticas em quaternions envolvem apenas produtos, sem funções trigonométricas, o que facilita a implementação embarcada em tempo real \cite{wie2008space}.
% Por essas razões, a formulação cinemática e dinâmica de atitude adotada neste trabalho utiliza quaternions unitários para a espaçonave perseguidora, evitando singularidades.

%%%%%%%%%%%%%%%%%%%%%%%%%
\subsubsection{Definição e componentes do quaternion}
\label{sec:componentes_quaternion}

O espaço quaternional é um espaço vetorial de quatro dimensões, composto por um elemento escalar e três elementos imaginários mutuamente ortogonais $\{i,j,k\}$, que obedecem às regras de multiplicação introduzidas por Hamilton (1844):
\begin{equation}
i^2 = j^2 = k^2 = ijk = -1,
\label{eq_hamilton_quadrados}
\end{equation}
\begin{equation}
ij = k,\quad jk = i,\quad ki = j,\quad
ji = -k,\quad kj = -i,\quad ik = -j.
\label{eq_hamilton_produtos}
\end{equation}

Um quaternion qualquer pode ser escrito como:
\begin{equation}
q = w + x\,i + y\,j + z\,k.
\label{eq_quaternion_forma}
\end{equation}

Ou, de forma vetorial compacta:

\begin{equation}
q = [\,w,\vec{v}\,] = [\,w,x,y,z\,].
\label{eq_quaternion_wxyz}
\end{equation}

Em que $w$ é a parte escalar e $\vec{v} = [x,y,z]$ é a parte vetorial.
Neste trabalho será adotada a notação $q = [w,x,y,z]^T$.
A magnitude (ou norma) de um quaternion é definida por:

\begin{equation}
\|q\| = \sqrt{w^2 + x^2 + y^2 + z^2}.
\label{eq_norma_quaternion}
\end{equation}

E o quaternion unitário é obtido normalizando-se $q$ de modo que $\|q\| = 1$.
Em aplicações de atitude, a condição de unidade constitui uma restrição do estado e pode ser empregada para monitorar o acúmulo de erros numéricos durante a integração.
Desvios de $\|q\|$ em relação a $1$ indicam erro numérico e podem ser corrigidos por renormalização do quaternion, evitando que a representação de atitude se degrade ao longo do tempo.

A relação entre os quaternions e a velocidade angular do corpo é obtida a partir da equação cinemática diferencial \cite{wie2008space}.
Considerando o vetor de estado quaternion:

\begin{equation}
\mathbf{q} =
\begin{bmatrix}
w \\
x \\
y \\
z
\end{bmatrix},
\qquad
\vec{\omega} =
\begin{bmatrix}
\omega_x \\ \omega_y \\ \omega_z
\end{bmatrix}.
\label{eq_q_omega_vetores}
\end{equation}

A equação cinemática pode ser escrita na forma matricial:

\begin{equation}
\begin{bmatrix}
\dot{w} \\[2pt]
\dot{x} \\[2pt]
\dot{y} \\[2pt]
\dot{z}
\end{bmatrix}
=
\frac{1}{2}
\begin{bmatrix}
0 & -\omega_x & -\omega_y & -\omega_z \\
\omega_x & 0 & \omega_z & -\omega_y \\
\omega_y & -\omega_z & 0 & \omega_x \\
\omega_z & \omega_y & -\omega_x & 0
\end{bmatrix}
\begin{bmatrix}
w \\ x \\ y \\ z
\end{bmatrix}.
\label{eq_cinematica_quaternion_wxyz}
\end{equation}

De forma compacta, pode-se escrever:

\begin{equation}
\dot{\mathbf{q}} = \tfrac{1}{2}\,\boldsymbol{\Omega}(\vec{\omega})\,\mathbf{q}.
\label{eq_quat_compacta}
\end{equation}

Em que $\boldsymbol{\Omega}(\vec{\omega})$ é a matriz $4\times 4$ definida em \eqref{eq_cinematica_quaternion_wxyz}.
